\documentclass[12pt]{article}

\usepackage{setspace}
\onehalfspace

%\usepackage[utf8]{inputenc}
%\usepackage{t1enc}
%\usepackage[T1]{fontenc}

\usepackage{amsmath}


%\DeclareMathOperator{\Exp}{Exp}
%\def\Exp{{\mathrm{Exp}}}

%\DeclareMathOperator{\Norm}{{\cal N}}
%\def\Norm{{\cal N}}

%\DeclareMathOperator{\Binom}{Binom}
%\def\Binom{{\mathrm{Binom}}}

%\DeclareMathOperator{\E}{E}
%\def\E{{\mathrm{E}}}

%\DeclareMathOperator{\D}{D}
%\def\D{{\mathrm{D}}}

%\renewcommand{\P}{\operatorname{P}}
%\def\P{{\mathrm{P}}}

%\DeclareMathOperator{\cov}{cov}
%\def\cov{{\mathrm{cov}}}


%\DeclareMathOperator{\corr}{corr}
%\def\corr{{\mathrm{corr}}}


%\renewcommand{\d}{\operatorname{d\! }} %integral \! a nyero
%\def\d{{\mathrm{d\! }}}


%\DeclareMathOperator{\R}{\mathbb{R}}
%\def\R{{\mathbb{R}}}

%\DeclareMathOperator{\Unif}{{\cal U}}
%\def\Unif{{\cal U}}

%\DeclareMathOperator{\Poi}{{Poisson}}
%\def\Poi{{\mathrm{Poisson}}}

%\def\ldots{{...\:}}


\usepackage{moodle}


\begin{document}

\begin{quiz}{gyakorló}

\begin{multi}[feedback={Egyenletes lineáris transzformáltja is egyenletes (a konstans elfajult esettől eltekintve)}]{impfvv5}
Ha $\xi\sim \Unif(0,1)$ akkor minden $a<b$-re $(b-a)\xi+a\sim \Unif(a,b)$.
\item* igaz
\item hamis
\end{multi}
\begin{multi}[feedback={Def.}]{impfvv8}
Ha $\xi\sim \Norm(0,1)$ akkor a sűrűségfüggvénye $f(x)=\frac{e^{-\frac{x^2}{2}}}{\sqrt{2\pi}}$.
\item* igaz
\item hamis
\end{multi}
\begin{multi}[feedback={$f\ge 0$ fontos (monoton nemcsökkenő $F$-nek a deriváltja)}]{sfv3}
Azt mondjuk, hogy egy $f:\R\to \R$ sűrűségfüggvény 
$$
\int_{-\infty}^{+\infty}f(x)\d x\ge 0
$$
\item igaz
\item* hamis
\end{multi}
\begin{multi}[feedback={Def.}]{val2}
A valószínűség additív, azaz $\P(A+B)=\P(A)+\P(B)$, ha $A$ és $B$ kizáróak.
\item* igaz
\item hamis
\end{multi}
\begin{multi}[feedback={Def.}]{flenvv1}
Azt mondjuk, hogy $\xi$ és $\eta$ diszkrét valségi változók függetlenek, ha
$$
\P(\xi=x,\eta=y)=\P(\xi=x)\P(\eta=y)
$$
bármely $x,y\in \R$-re teljesül.
\item* igaz
\item hamis
\end{multi}
\begin{multi}[feedback={de-Morgan azonosság}]{esem1}
Az $\overline{A+B}$ és $\overline{A} \cdot \overline{B}$ események megegyeznek.
\item* igaz
\item hamis
\end{multi}
\begin{multi}[feedback={Definíció.}]{minta6}
Egy $X_1,\ldots,X_n\sim X$ minta korrigált empirikus szórásnégyzetén a 
$$
s^2=\frac{\sum_{k} (X_k-\overline{X})^2}{n-1}
$$
mennyiséget értjük.
\item* igaz
\item hamis
\end{multi}
\begin{multi}[feedback={A mintaátlag torzítatlan becslése a sokasági (elméleti) várható értéknek (átlagnak)}]{minta4}
Egy $X_1,\ldots,X_n\sim X$ minta esetén $\E(\overline{X})=\E(X)$.
\item* igaz
\item hamis
\end{multi}
\begin{multi}[feedback={$A=B$?}]{val3}
A valószínűség monoton, azaz $\P(A)<\P(B)$, ha $A\subseteq B$.
\item igaz
\item* hamis
\end{multi}
\begin{multi}[feedback={Csak ha az elemi események egyforma valségűek.}]{val5}
Bármely esemény valsége kiszámolható a $\frac{\text{kedvező}}{\text{összes}}$ képlettel.
\item igaz
\item* hamis
\end{multi}
\begin{multi}[feedback={Def.}]{impdvv6}
Azt mondjuk, hogy $\xi\sim \Binom(n,p)$, ha 
$$
\P(\xi=k)=\binom{n}{k}p^k(1-p)\ \ \ (k=0,1\ldots n)
$$
valamilyen $n>0$ és $p\in[0,1]$ esetén.
\item igaz
\item* hamis
\end{multi}
\begin{multi}[feedback={wiki}]{impfvv1}
Egy $\xi\sim \Unif(a,b)$ várható értéke: $\frac{a+b}{2}$
\item* igaz
\item hamis
\end{multi}
\begin{multi}[feedback={Def.}]{dvv6}
Egy $\xi$ valségi változó az $x_1,\ldots, x_n$ értékeket veheti fel. Ekkor a szórása:
$$
\D(\xi)=\sum_{k=1}^n \P(\xi=x_k)(x_k-\E(\xi))^2
$$
\item igaz
\item* hamis
\end{multi}
\begin{multi}[feedback={Ezek az $t$-próba használatának előzményei.}]{proba3}
Ha a minta normális sokaságból származik, nem ismert a szórás és a sokasági várható értékre (átlagra) vonatkozó kérdésről 
akarunk dönteni akkor $t$-próbát használhatunk.
\item* igaz
\item hamis
\end{multi}
\begin{multi}[feedback={A fenti formula a $u$-próba ún. próbastatisztikája.}]{proba8}
Az $t$-próbánál a 
$$
u=\frac{\overline{X}-\mu_0}{\sigma}\sqrt{n}.
$$
mennyiség segítségével döntünk a hipotéziseinkről.
\item igaz
\item* hamis
\end{multi}
\begin{multi}[feedback={Definíció.}]{esem8}
Az $A$ és $B$ események pont akkor függetlenek, ha $\P(A+B)=\P(A)+\P(B)$.
\item igaz
\item* hamis
\end{multi}
\begin{multi}[feedback={A páronkénti gyengébb tulajdonság.}]{felt8}
A páronkénti függetlenségből következik a teljes.
\item igaz
\item* hamis
\end{multi}
\begin{multi}[feedback={Bayes-tétel}]{felt6}
Ha $A_1,\ldots,A_n$ teljes eseményrendszer pozitív valségű tagokkal és $\P(B)>0$ akkor minden $m$-re:
$$
\P(A_m|B)=\frac{\P(B|A_m)\P(A_m)}{\sum_{k=1}^n \P(B|A_k)\P(A_k)}
$$
\item* igaz
\item hamis
\end{multi}
\begin{multi}[feedback={A fenti formula az $u$-próba ún. próbastatisztikája.}]{proba5}
Az $u$-próbánál a 
$$
u=\frac{\overline{X}-\mu_0}{\sigma}\sqrt{n}
$$
mennyiség segítségével döntünk a hipotéziseinkről.
\item* igaz
\item hamis
\end{multi}
\begin{multi}[feedback={Def.}]{flenvv4}
Ha $\xi$ és $\eta$ diszkrét valségi változóknak létezik a szórása és nemnulla, akkor a 
korrelációjuk
$$
\corr(\xi,\eta)=\frac{\cov(\xi,\eta)}{\D(\xi)\D(\eta)}
$$
\item* igaz
\item hamis
\end{multi}
\begin{multi}[feedback={$\xi=\xi_1+\ldots+\xi_n$, ahol $\xi_k$ Bernoulli és $\E(\xi_k)=p$+$\E$-additívitása}]{impdvv2}
Egy $\xi\sim \Binom(n,p)$ várható értéke: $n+p$
\item igaz
\item* hamis
\end{multi}
\begin{multi}[feedback={Def.}]{sfv1}
Azt mondjuk, hogy egy $f:\R\to \R$ sűrűségfüggvény, ha 
$$
f\ge 0 \ \ \text{és}\ \ \int_{-\infty}^{+\infty}f(x)\d x=1
$$
\item* igaz
\item hamis
\end{multi}
\begin{multi}[feedback={Ebből nem jön ki a pl. a monotonitás vagy a balról folytonosság.}]{eofv2}
Azt mondjuk, hogy egy $F:\R\to \R$ eloszlásfüggvény, ha $F\ge 0$ és 
$$
\int_{-\infty}^{+\infty}F(x)\d x=1
$$
\item igaz
\item* hamis
\end{multi}
\begin{multi}[feedback={Tulajdonság.}]{flenvv6}
Ha $\xi$ és $\eta$ valségi változóknak létezik a szórása, akkor a 
kovarianciájuk
$$
\cov(\xi,\eta)=\E(\xi\eta)-\E(\xi)\E(\eta)
$$
\item* igaz
\item hamis
\end{multi}
\begin{multi}[feedback={Def.}]{dvv9}
Egy $\xi$ valségi változó az $x_1,\ldots, x_n$ értékeket veheti fel és $p_k=\P(\xi=x_k)$. 
Ekkor a második momentuma:
$$
\E(\xi^2)=\sum_{k=1}^n x_k p_k^2
$$
\item igaz
\item* hamis
\end{multi}
\begin{multi}[feedback={Hiszen $F(x)=\P(\xi<x)$, és ez mindig értelmes.}]{eofv4}
Minden valségi változónak van eloszlásfüggvénye.
\item* igaz
\item hamis
\end{multi}
\begin{multi}[feedback={Def.}]{dvv7}
Egy $\xi$ valségi változó az $x_1,\ldots, x_n$ értékeket veheti fel. Ekkor a szórásnégyzete:
$$
\D^2(\xi)=\sum_{k=1}^n \P(\xi=x_k)(x_k-\E(\xi))^2
$$
\item* igaz
\item hamis
\end{multi}
\begin{multi}[feedback={Def.}]{sfv9}
Ha $\xi$-nek $f$ a sűrűségfüggvénye és $F$ az eloszlásfüggvénye, akkor bármely $b$-re:
$$
F(b)=\int_{\infty}^{b}f(x)\d x
$$
\item* igaz
\item hamis
\end{multi}
\begin{multi}[feedback={wiki}]{impfvv4}
Egy $\xi\sim \Unif(a,b)$ szórásnégyzete: $\frac{(a+b)^2}{2}$
\item igaz
\item* hamis
\end{multi}
\begin{multi}[feedback={Például: legyen a kísérlet random húzás a $[0,1]$-ből és $A=\{\frac{1}{2}\}$}]{esem9}
Ha $\P(A)=0$ akkor $A$ bekövetkezhet.
\item* igaz
\item hamis
\end{multi}
\begin{multi}[feedback={Def.+$P(AB)=P(A)$ a feltételekből.}]{felt1}
$A$-nak a $B$-re ($\P(B)>0$) vonatkozó feltételes valsége $\P(A|B)=\frac{\P(A)}{\P(B)}$, ha $A$ maga után vonja $B$-t.
\item* igaz
\item hamis
\end{multi}
\begin{multi}[feedback={Def.}]{sfv2}
Egy $f:\R\to [0,1]$ pontosan akkor sűrűségfüggvény, ha monoton nemcsökkenő, 
balról folytonos, $\lim_{x\to+\infty}f(x)=1$ és $\lim_{x\to-\infty}f(x)=0$
\item igaz
\item* hamis
\end{multi}
\begin{multi}[feedback={Ilyesmi a $\corr$-ra igaz.}]{flenvv8}
Ha a $\xi$ és $\eta$ valségi változóknak létezik a kovarianciája, akkor $-1\le \cov(\xi,\eta)\le 1$.
\item igaz
\item* hamis
\end{multi}
\begin{multi}[feedback={de-Morgan azonosság}]{esem2}
Az $\overline{A+B}$ és $\overline{A} \cdot \overline{B}$ események kizárják egymást.
\item igaz
\item* hamis
\end{multi}
\begin{multi}[feedback={Páronként független valségi változókra a szórásnégyzet additív.}]{dvv13}
$\xi$ és $\eta$ független, véges értékkészletű valségi változókra:
$$
\D^2(\xi+\eta)=\D^2(\xi)+\D^2(\eta)
$$
\item* igaz
\item hamis
\end{multi}
\begin{multi}[feedback={A feltételek mellett a fenti mennyiség sztenderd normális, ha $H_0$ igaz.}]{proba10}
Ha egy $16$ elemű minta normális sokaságból származik, $\sigma=3$ és hipotéziseink:
$$
H_0: \mu=10
$$
$$
H_1: \mu\neq 10
$$
akkor az 
$$
\frac{\overline{X}-10}{3}4
$$
valségi változó $15$-szabadsági fokú Student eloszlású ha $H_0$ igaz.
\item igaz
\item* hamis
\end{multi}
\begin{multi}[feedback={Az adott feltételek mellet a várható érték egy súlyozott közép.}]{dvv5}
Egy $\xi$ valségi változó az $x_1,\ldots, x_n$ értékeket veheti fel. Ekkor:
$$
\min_k(x_k)\le \E(\xi) \le \max_k(x_k)
$$
\item* igaz
\item hamis
\end{multi}
\begin{multi}[feedback={Def.}]{impdvv5}
Azt mondjuk, hogy $\xi\sim \Binom(n,p)$, ha 
$$
\P(\xi=k)=\binom{n}{k}p^k(1-p)^{n-k}\ \ \ k=0\ldots n
$$
valamilyen $n>0$ és $p\in[0,1]$ esetén.
\item* igaz
\item hamis
\end{multi}
\begin{multi}[feedback={A balról folytonosság kell.}]{eofv3}
Egy $F:\R\to [0,1]$ pontosan akkor eloszlásfüggvény, ha monoton nemcsökkenő, 
$\lim_{x\to+\infty}F(x)=1$ és $\lim_{x\to-\infty}F(x)=0$
\item igaz
\item* hamis
\end{multi}
\begin{multi}[feedback={Páronként független valségi változókra a szórásnégyzet additív, általában nem.}]{dvv14}
A $\xi$ és $\eta$ valségi változókra:
$$
\D^2(\xi+\eta)=\D^2(\xi)+\D^2(\eta)
$$
feltéve, hogy létezik a szórásuk.
\item igaz
\item* hamis
\end{multi}
\begin{multi}[feedback={Az $s^2$ torzítatlan becslése a sokasági (elméleti) szórásnégyzetnek}]{minta8}
Az $s^2$ korrigált empirikus szórásnégyzetre $\E(s^2)=\E(X)$.
\item igaz
\item* hamis
\end{multi}
\begin{multi}[feedback={Komplementer események.}]{eofv9}
Bármely $\xi$-re és $a$-ra $F_{\xi}(a)+\P(\xi\ge a)=1$.
\item* igaz
\item hamis
\end{multi}
\begin{multi}[feedback={A feltételek mellett a fenti Student (t) eloszlású.}]{proba12}
Ha egy $16$ elemű minta normális sokaságból származik, a szórás nem ismert, $s=3$ és hipotéziseink:
$$
H_0: \mu=10 \\
$$
$$
H_1: \mu\neq 10
$$
akkor az 
$$
\frac{\overline{X}-10}{3}4
$$
valségi változó sztenderd normális eloszlású ha $H_0$ igaz.
\item igaz
\item* hamis
\end{multi}
\begin{multi}[feedback={Def.}]{impfvv11}
Ha $\xi\sim \Unif(a,b)$ akkor a sűrűségfüggvénye 
$$
f(x)=
\begin{cases}
\frac{1}{b-a} & \ \ a<x<b \\
0 & \text{máskor} 
\end{cases}
$$
\item* igaz
\item hamis
\end{multi}
\end{quiz}

\end{document}

